% 第12章（本部分第5章）图：贯穿例 A=[[4,1],[1,3]], b=(1,2), x0=0 上 CG 与最速下降的路径。
% 坐标、等值线半轴与转角由同目录 data.py 用有理数精确计算后写入 path-data.tex。
% CG：x1=(1/4,1/2), x2=x_star=(1/11,7/11)；最速下降：x1 相同，x2=(1/12,7/12), x3=(5/48,5/8), x4=(13/144,91/144)。
% 左：总览（60 mm/单位）；右：x_star 附近放大（600 mm/单位，放大10倍）。
% 由 data.py 生成，勿手改。单位 mm；O=总览（60 mm/单位，原点 x0），Z=放大（600 mm/单位，原点 x_star）。
\def\CGangle{31.7175}
\def\Oxs{5.455}\def\Oys{38.182}
\def\Olevels{32.604/45.398,9.412/13.105,2.717/3.783,0.784/1.092,0.226/0.315,0.065/0.091}
\def\Ocg{(0.000,0.000) -- (15.000,30.000) -- (5.455,38.182)}
\def\Osd{(0.000,0.000) -- (15.000,30.000) -- (5.000,35.000) -- (6.250,37.500) -- (5.417,37.917) -- (5.521,38.125) -- (5.451,38.160)}
\def\Osdpts{5.000/35.000,6.250/37.500,5.417/37.917}
\def\Zxs{0.000}\def\Zys{0.000}
\def\Zlevels{326.041/453.976,94.120/131.052,27.170/37.831,7.843/10.921,2.264/3.153,0.654/0.910}
\def\Zcg{(-54.545,-381.818) -- (95.455,-81.818) -- (0.000,0.000)}
\def\Zsd{(-54.545,-381.818) -- (95.455,-81.818) -- (-4.545,-31.818) -- (7.955,-6.818) -- (-0.379,-2.652) -- (0.663,-0.568) -- (-0.032,-0.221)}
\def\Zsdpts{-4.545/-31.818,7.955/-6.818,-0.379/-2.652}

\begin{tikzpicture}[x=1mm,y=1mm,>=Stealth,
  note/.style={font=\FigureNoteFont\fontsize{7.5}{9}\selectfont,text=FigureMuted},
  pt/.style={circle,inner sep=0pt,minimum size=1.5mm,draw=FigureStroke,line width=.5pt},
  spt/.style={circle,inner sep=0pt,minimum size=1.1mm,draw=FigureYellow,fill=white,line width=.6pt}]
% ---------- 左：总览 ----------
\def\Lx{-7}\def\Rx{24}\def\By{-5}\def\Ty{44}
\begin{scope}
\clip (\Lx,\By) rectangle (\Rx,\Ty);
\foreach \a/\b in \Olevels{
  \draw[FigureBlue,line width=.5pt,rotate around={\CGangle:(\Oxs,\Oys)}]
    (\Oxs,\Oys) ellipse [x radius=\a, y radius=\b];}
\end{scope}
\draw[FigureGrid,line width=.4pt] (\Lx,\By) rectangle (\Rx,\Ty);
\draw[FigureGreen,line width=1.1pt,->] \Ocg;
\node[pt,fill=white] at (0,0) {};
\node[pt,fill=white] at (15,30) {};
\node[pt,fill=FigureGreen] at (\Oxs,\Oys) {};
\node[anchor=west] at (1.2,0) {$\bm x_0$};
\node[anchor=west] at (16,30) {$\bm x_1$};
\node[anchor=east] at ({\Oxs-1.6},\Oys) {$\bm x_\star$};
\draw[FigureStroke,line width=.5pt] ({\Oxs-1.4},{\Oys-3.8}) rectangle ({\Oxs+1.4},{\Oys+0.6});
\node[font=\figurefont\footnotesize] at ({(\Lx+\Rx)/2},{\Ty+3.5}) {总览};
% ---------- 右：放大 ----------
\begin{scope}[shift={(48,35)}]
\def\zL{-14}\def\zR{14}\def\zB{-38}\def\zT{6}
\begin{scope}
\clip (\zL,\zB) rectangle (\zR,\zT);
\foreach \a/\b in \Zlevels{
  \draw[FigureBlue,line width=.5pt,rotate around={\CGangle:(0,0)}]
    (0,0) ellipse [x radius=\a, y radius=\b];}
\draw[FigureYellow,line width=.9pt,dashed] \Zsd;
\draw[FigureGreen,line width=1.1pt,->] \Zcg;
\foreach \u/\v in \Zsdpts{\node[spt] at (\u,\v) {};}
\end{scope}
\draw[FigureStroke,line width=.5pt] (\zL,\zB) rectangle (\zR,\zT);
\node[pt,fill=FigureGreen] at (0,0) {};
\node[anchor=south east] at (-0.6,0.4) {$\bm x_\star$};
\node[anchor=east,text=FigureYellow] at (-5.4,-31.8) {$\bm x_2$};
\node[anchor=north,text=FigureYellow] at (8.6,-7.6) {$\bm x_3$};
\node[anchor=east,text=FigureYellow] at (-1.2,-2.9) {$\bm x_4$};
\node[font=\figurefont\footnotesize] at (0,{\zT+3.5}) {$\bm x_\star$附近放大10倍};
\end{scope}
\draw[FigureStroke,line width=.4pt,densely dotted] ({\Oxs+1.4},{\Oys+0.6}) -- (34,41);
\draw[FigureStroke,line width=.4pt,densely dotted] ({\Oxs+1.4},{\Oys-3.8}) -- (34,-3);
% ---------- 图例 ----------
\begin{scope}[shift={(-5,-11)}]
\draw[FigureGreen,line width=1.1pt] (0,0)--(6,0); \node[note,anchor=west] at (7,0) {CG};
\draw[FigureYellow,line width=.9pt,dashed] (20,0)--(26,0); \node[note,anchor=west] at (27,0) {最速下降};
\draw[FigureBlue,line width=.5pt] (47,0)--(53,0); \node[note,anchor=west] at (54,0) {$\phi$的等值线};
\end{scope}
\end{tikzpicture}
